This course is organised as follows. Lectures take place on Mondays and lab sessions on Tuesdays. During the lectures, key concepts are introduced from a conceptual and theoretical perspective, accompanied by examples of Python implementations.

Lab sessions begin with knowledge questions about the week’s literature and lecture. Students then work on the assigned exercises and continue working on the group assignment once these are completed. The tutorial lecturer (Roeland) will circulate to help with questions. If several students encounter the same issue, it will be discussed plenarily. Active participation in the tutorial meetings is important to ensure that students understand the material and can work effectively on the assignments.

\section*{Week 1: Course introduction \& Working with textual data}

\subsection*{Monday, 30--3: Lecture (Anne)}

In this first lecture, we introduce the course and its structure. We also explore how computational methods can be used in communication science by examining what we can learn from analysing textual data.

\subsection*{Tuesday, 31--3: Lab session}

\textsc{\ding{52} Read this manual and inspect the course Canvas page.}\\
\textsc{\ding{52} Ensure that your computer is ready to use for the course.}\\
\textsc{\ding{52} Read in advance: \cite{vanAtteveldt2018}.}\\
\textsc{\ding{52} Read in advance: \cite{Boumans2016}.}\\

During the first tutorial meeting we introduce the course goals and policies and apply the lecture material to a first text analysis exercise. We also revisit some basic concepts from CCS-1 and practise simple top-down and bottom-up approaches to text analysis.

Students will form groups of approximately four students for the group assignment. If you complete the in-class exercises early, you may begin exploring the dataset, inspecting relevant variables, and starting the data cleaning process.

\section*{Week 2: Bottom-up approaches to text analysis}

\subsection*{Monday, 6--4: Easter Monday --- no lecture}

\subsection*{Tuesday, 7--4: Lecture (Anne)}

\emph{This week the lecture takes place on Tuesday instead of Monday due to Easter Monday. There will be no separate lab session.}

\textsc{\ding{52} Read in advance: Chapter 10 “Text as data” \cite{van_atteveldt_computational_2022}.}\\
\textsc{\ding{52} Read in advance: \cite{kroon2024advancing}. Read from “From BoW to Large Language Models: Revolutions in text representations” until “The rise of the transformer” (pp. 147–149).}

In this lecture we discuss different ways to represent text data, including Bag of Words, TF-IDF, and word embeddings.

Students will practise these concepts independently by transforming textual data into matrices using \texttt{sklearn}'s vectorizers, measuring similarity between documents, and experimenting with a simple LDA model.

\section*{Week 3: Recommender systems}

\subsection*{Monday, 13--4: Lecture (Anne)} 

\textsc{\ding{52} Read in advance: \cite{Moller2018}.}\\
\textsc{\ding{52} Read in advance: \cite{Loecherbach2020}.}\\

This lecture introduces recommender systems. We discuss different types of recommendation approaches and how they rely on earlier concepts such as preprocessing text data, transforming texts into matrices, and calculating similarity between documents.

\subsection*{Tuesday, 14--4: Lab session}

During this lab session you will complete the first set of multiple-choice questions and practise building a simple recommender system.

\section*{Week 4: Introduction to supervised machine learning \\ Group presentation recommender system}

\subsection*{Monday, 20--4: Lecture (Saurabh)}

\textsc{\ding{52} Read in advance: \cite{vermeer_seeing_2019}.}\\
\textsc{\ding{52} Read in advance: \cite{meppelink_reliable_2021}.}\\

This lecture introduces supervised machine learning and discusses its basic principles using practical examples and code.

\subsection*{Tuesday, 21--4: Lab session}

During this lab session groups will present their recommender system assignment.

\subsection*{Deadline group project: Friday 24--4 at 23:59}

The deadline for submitting the group assignment is \textbf{Friday 24--4 at 23:59}.

\section*{Week 5: Teaching-free week}

\subsection*{Monday, 27--4: No lecture --- UvA teaching-free week}
\subsection*{Tuesday, 28--4: No lab session --- UvA teaching-free week}

\textsc{\ding{52} Take a break}

\section*{Week 6}

\subsection*{Monday, 4--5: No lecture --- Collectieve roostervrije dag}
\subsection*{Tuesday, 5--5: Liberation Day --- no lab session}

\section*{Week 7: Validation in supervised machine learning}

\subsection*{Monday, 11--5: Lecture (Saurabh)}

\textsc{\ding{52} Read in advance: \cite{birkenmaier_search_2023}.}\\

In this lecture we discuss validation techniques for supervised machine learning models.

\subsection*{Tuesday, 12--5: Lab session}

During the lab session we practise evaluating and validating machine learning models.

\section*{Week 8: Coding in an academic context}

\subsection*{Monday, 18--5: Lecture (Saurabh)}

\textsc{\ding{52} Read in advance: \cite{hube_understanding_2019}.}\\
\textsc{\ding{52} Read in advance: \cite{bender_dangers_2021}.}\\

In this lecture we reflect on the responsible use of computational methods and machine learning in academic research.

\subsection*{Tuesday, 19--5: Lab session}

The final lab session focuses on responsible research coding and reproducible research practices.

\section*{Week 9: Exam week}

\subsection*{Week of 1 June: Exam inspection}

\subsection*{Week of 8 June: Resit}

\subsection*{Week of 15 June: Resit inspection}